\section{Transformer for fMRI：从 ROI 序列到脑网络表示}

Karan 的案例展示了“同一架构、不同数据几何”。fMRI 不是图片分类：输入是多个 brain regions of interest（ROI）随时间变化的低频信号，噪声高、样本少且个体差异大。成功迁移需要重新定义 token、mask objective、cross-attention 与 downstream validation。

\subsection{数据接口：ROI-by-time 序列}

fMRI 通过 blood-oxygen-level-dependent signal 间接反映神经活动。将 $R$ 个 ROI、$T$ 个时间点写成 $X\in\mathbb{R}^{T\times R}$ 后，可把时间窗或 ROI embedding 当 token。相关矩阵能描述静态 functional connectivity，却难以表示动态、非线性和方向依赖。


\lecturefigure{slide-099.jpg}{fMRI：functional connectivity、averaging 与 correlation maps}{CS25 V5 official deck, p.~99}


\lecturefigure{slide-100.jpg}{Brain functional networks：ROI 分区与网络标签}{CS25 V5 official deck, p.~100}


\lecturefigure{slide-101.jpg}{Early ML：相关矩阵、手工特征与传统分类器}{CS25 V5 official deck, p.~101}

\begin{warningbox}{脑区相关不等于因果连接}
相关矩阵受共同输入、预处理和采样频率影响。Transformer 学到的 attention 或 embedding 也不能直接解释为神经因果机制；它们首先是预测表示，需要与 intervention、跨被试复现和临床 outcome 分开验证。
\end{warningbox}

\subsection{Foundation model objective：masked prediction}

课程采用 data-driven self-supervision：遮住部分 ROI/time tokens，用其余脑活动预测缺失区域。设 mask 集为 $M$，可写为
\[
\mathcal{L}_{\text{mask}}=\sum_{(t,r)\in M}\|x_{t,r}-\hat{x}_{t,r}\|_2^2.
\]
该目标迫使模型学习跨区域与跨时间依赖，而不是依赖人工定义的疾病标签；但预训练收益取决于 cohort diversity、scanner/site shift 和 mask design。


\lecturefigure{slide-102.jpg}{Foundation models for fMRI：masked modeling 与可迁移 encoder}{CS25 V5 official deck, p.~102}


\lecturefigure{slide-103.jpg}{Data-driven approach：从全脑序列自动学习表示}{CS25 V5 official deck, p.~103}

\teachervoice{00:41:34--00:45:16，Karan 把 ROI-by-time 数据直接输入 Transformer，并通过遮挡区域重建构造 self-supervised objective。课堂洞见是：先找一个能利用无标签数据的任务，再讨论 downstream diagnosis。}

\subsection{Cross-attention 与架构变体}

若给定 unmasked context $C$ 和 target region queries $Q_T$，cross-attention 为
\[
H_T=\operatorname{softmax}\!\left(\frac{Q_TK_C^{\top}}{\sqrt{d_k}}\right)V_C.
\]
目标脑区主动查询其余网络，而不是简单把所有 ROI 平均。模型还可在 region、time 和 subject 维度采用不同 encoder/decoder 结构；选择应由泛化和 compute 验证，而不是照搬 NLP block。


\lecturefigure{slide-104.jpg}{Cross-attention：用可见脑区预测被遮挡 target region}{CS25 V5 official deck, p.~104}


\lecturefigure{slide-105.jpg}{fMRI model architectures：encoder-decoder 与 attention 连接}{CS25 V5 official deck, p.~105}

\begin{knowledgebox}{深读：fMRI 架构选择要服从数据轴而不是 NLP 惯例}
脑信号同时有 time、region、subject 和 site 四个主要变化轴。若 token 表示时间窗，attention 擅长学习跨时间依赖；若 token 表示 ROI，则可学习网络间 interaction；若把二者直接展平，序列会很长且语义混杂。Cross-attention 可以让 target region 查询 context regions，encoder-decoder 可以隔离观测与预测，但这些结构都需要与简单 temporal convolution、graph model 和 correlation baseline 比较。Mask ratio 过低会让重建太容易，过高则丢失可预测上下文。Subject embedding 可能吸收个体差异，也可能泄露身份。最关键的 split 是按 subject/site 切分而非随机 time points，否则临近扫描的自相关会虚高结果。架构图只是 hypothesis space，严谨 protocol 才决定 evidence quality。
\end{knowledgebox}

\subsection{结果：先看 representation，再看 downstream task}

结果页同时展示 reconstruction/modeling 指标、network dependence 与 downstream classification。读图应先检查 train/test split 是否跨 subject/site，再比较 correlation baseline、linear model 与 foundation representation；若只在同站点随机切分，模型可能利用 scanner 或人口学 shortcut。


\lecturefigure{slide-106.jpg}{Modeling results I：重建指标与脑网络分组}{CS25 V5 official deck, p.~106}


\lecturefigure{slide-107.jpg}{Modeling results II：不同架构与区域的比较}{CS25 V5 official deck, p.~107}


\lecturefigure{slide-108.jpg}{Network dependence：不同脑网络间的预测依赖}{CS25 V5 official deck, p.~108}


\lecturefigure{slide-109.jpg}{Downstream task：embedding 用于疾病或个体表征}{CS25 V5 official deck, p.~109}


\lecturefigure{slide-110.jpg}{Downstream results：foundation representation 与传统 baseline}{CS25 V5 official deck, p.~110}

\begin{knowledgebox}{读图：fMRI 结果页的证据顺序}
先看 representation 是否在 held-out signal reconstruction 上有效，再看不同 network 的依赖结构，最后才看 downstream disease classification。Slide 110 若给出更高 accuracy，仍要追问 subject-level split、site leakage、class balance 和 calibration。饼图或 embedding 可视化只能描述模型如何分组，不能证明临床因果。最强结论是“该表示在这套 protocol 下提供增量预测信息”，而不是“模型理解了大脑”或可直接用于诊断。
\end{knowledgebox}

\teachervoice{00:45:16--00:48:02，讲者把 cross-attention 解释为用全脑上下文预测目标区域，并报告 learned representation 在 downstream disease task 上优于 correlation/linear baselines。讲义提醒：这是特定数据和 split 的证据，不是临床部署结论。}

\begin{knowledgebox}{神经科学 foundation model 的验证层级}
\begin{enumerate}
  \item masked reconstruction 是否优于简单 baseline；
  \item representation 是否跨 subject/site 稳定；
  \item downstream task 是否在外部 cohort 复现；
  \item 是否具有 calibration、公平性与临床 utility；
  \item 模型解释是否与已知神经机制和 intervention 一致。
\end{enumerate}
前两层通过不代表最后三层自动成立。
\end{knowledgebox}

\subsection{本章小结}

fMRI 案例说明 Transformer 的可迁移性来自 token interaction，而不是 NLP 名称。ROI/time tokenization、masked objective、cross-attention 与跨 cohort 验证共同决定成败；模型可学习有用表示，但 attention 和相关性都不能直接升级为脑因果解释。

